\documentclass[12pt,a4paper]{article}

\usepackage[romanian]{babel}
\usepackage[T1]{fontenc}
\usepackage[utf8]{inputenc}
\usepackage{lmodern}
\usepackage{enumitem}

\title{Modelarea Conceptuală a Bazelor de Date \\ ERD, Entități, Relații, Atribute}
\author{}
\date{}

\begin{document}

\maketitle

\section{Entități}
O \textbf{entitate} reprezintă un obiect sau concept din lumea reală care poate fi
identificat în mod unic.

\subsection*{Entități Strong}
O entitate \textbf{strong} (puternică):
\begin{itemize}[noitemsep]
    \item are existență independentă;
    \item are o cheie primară proprie;
    \item nu depinde de alte entități pentru identificare.
\end{itemize}

\textbf{Exemplu:} \texttt{Student(StudID, Nume, Email)}

\subsection*{Entități Weak}
O entitate \textbf{weak} (slabă):
\begin{itemize}[noitemsep]
    \item nu poate fi identificată fără o entitate strong;
    \item depinde de cheia primară a entității strong;
    \item are o cheie parțială (\textit{discriminator}).
\end{itemize}

\textbf{Exemplu:} \texttt{TelefonStudent(StudID, NrTelefon)}  
Aici \texttt{StudID} provine din entitatea strong \texttt{Student}.

\section{Atribute}
Atributele descriu proprietățile unei entități.

\subsection*{Tipuri de atribute}
\begin{itemize}[noitemsep]
    \item \textbf{Simple}: nu pot fi divizate (ex.: \texttt{Nume}).
    \item \textbf{Compuse}: pot fi împărțite în sub-atribute (ex.: \texttt{Adresă}).
    \item \textbf{Multivaloare}: pot avea mai multe valori (ex.: \texttt{Telefon}).
    \item \textbf{Derivate}: se calculează din alte atribute (ex.: \texttt{Vârstă}).
\end{itemize}

\section{Chei}
\subsection*{Cheia Primară}
\textbf{Cheia primară} este un atribut sau un set de atribute care identifică în mod
unic fiecare tuplu dintr-o entitate.

Proprietăți:
\begin{itemize}[noitemsep]
    \item unicitate;
    \item minimalitate;
    \item non-null.
\end{itemize}

\subsection*{Chei Candidate și Superchei}
\begin{itemize}[noitemsep]
    \item \textbf{Chei candidate}: toate seturile minimale de atribute care pot fi chei primare.
    \item \textbf{Superchei}: orice set de atribute care identifică unic tuplurile, dar nu neapărat minimal.
\end{itemize}

\section{Relații}
O \textbf{relație} descrie asocierea dintre două sau mai multe entități. Relațiile se clasifică după aritate și cardinalitate.

\subsection*{Relații Unare (Recursive)}
O relație este \textbf{unară} dacă o entitate se relaționează cu ea însăși.

\textbf{Exemple:}
\begin{itemize}[noitemsep]
    \item \textbf{Angajat -- seSubordonează -- Angajat} (relație ierarhică)
    \item \textbf{Produs -- esteComponență -- Produs} (structuri compuse)
\end{itemize}

\subsection*{Relații Binare}
Relațiile \textbf{binare} implică două entități distincte și sunt cele mai frecvente.

\textbf{Exemple:}
\begin{itemize}[noitemsep]
    \item Student -- seÎnscrie -- Curs
    \item Profesor -- predă -- Curs
    \item Student -- areCard -- CardAcces
\end{itemize}

\subsection*{Relații N-are (Ternare și Superioare)}
O relație este \textbf{ternară} dacă implică trei entități, sau \textbf{n-ară} dacă implică mai multe.

\textbf{Exemplu ternar:}
\begin{itemize}[noitemsep]
    \item \textbf{Furnizor -- livrează -- Produs -- către -- Magazin}
\end{itemize}

\section{Cardinalitatea Relațiilor}
Cardinalitatea definește numărul minim și maxim de participări ale unei entități
într-o relație.

\subsection*{Tipuri de cardinalități}
\begin{itemize}[noitemsep]
    \item \textbf{(0,1)} participare opțională, maxim un element;
    \item \textbf{(1,1)} participare obligatorie, exact un element;
    \item \textbf{(0,N)} participare opțională, mai multe elemente;
    \item \textbf{(1,N)} participare obligatorie, mai multe elemente.
\end{itemize}

\subsection*{Tipuri de relații după cardinalitate}
\begin{itemize}[noitemsep]
    \item \textbf{1:1}: un student are un singur card de acces.
    \item \textbf{1:N}: un profesor predă mai multe cursuri.
    \item \textbf{M:N}: un student se înscrie la mai multe cursuri, iar un curs are mai mulți studenți.
\end{itemize}

\section{Mini-ERD Textual Complet}

\subsection*{Entități}
\begin{itemize}[noitemsep]
    \item \textbf{Student(StudID, Nume, Email, DataNasterii)} -- entitate strong
    \item \textbf{Curs(CursID, Denumire, Credite)} -- entitate strong
    \item \textbf{Profesor(ProfID, Nume, Titlu)} -- entitate strong
    \item \textbf{TelefonStudent(StudID, NrTelefon)} -- entitate weak
\end{itemize}

\subsection*{Atribute speciale}
\begin{itemize}[noitemsep]
    \item \texttt{Adresă} pentru Student poate fi atribut compus: (Stradă, Număr, Oraș).
    \item \texttt{NrTelefon} este atribut multivaloare pentru Student, modelat prin entitatea weak \texttt{TelefonStudent}.
    \item \texttt{Vârstă} poate fi atribut derivat din \texttt{DataNasterii}.
\end{itemize}

\subsection*{Relații și Cardinalități}
\begin{itemize}[noitemsep]
    \item \textbf{Student -- TelefonStudent}: relație 1:N (un student poate avea mai multe telefoane).
    \item \textbf{Student -- Curs}: relație M:N prin entitatea asociativă:
    \begin{itemize}[noitemsep]
        \item \textbf{Inscriere(StudID, CursID, Nota)}
    \end{itemize}
    \item \textbf{Profesor -- Curs}: relație 1:N (un profesor poate preda mai multe cursuri).
    \item \textbf{Angajat -- seSubordonează -- Angajat}: relație unară 1:N (un manager coordonează mai mulți angajați).
\end{itemize}

\subsection*{Chei Primare în Mini-ERD}
\begin{itemize}[noitemsep]
    \item \textbf{Student}: \texttt{StudID} (cheie primară).
    \item \textbf{Curs}: \texttt{CursID}.
    \item \textbf{Profesor}: \texttt{ProfID}.
    \item \textbf{TelefonStudent}: \texttt{(StudID, NrTelefon)} (cheie compusă).
    \item \textbf{Inscriere}: \texttt{(StudID, CursID)} (cheie compusă).
\end{itemize}

\end{document}
